%\documentclass[11pt]{amsart}
\documentclass[12pt,letterpaper]{report}
% Font and Encoding
\usepackage[T1]{fontenc}
\usepackage{times}  % Times Roman font
% Geometry for page layout
\usepackage[letterpaper,top=1in, bottom=1.25in, left=1.25in, right=1in, headheight=0.5in, headsep=0.25in]{geometry}
% Graphics and Images
\usepackage{graphicx}
\usepackage{epstopdf}
\DeclareGraphicsRule{.tif}{png}{.png}{`convert #1 `dirname #1`/`basename #1 .tif`.png}
% Math Packages
\usepackage{amsmath}
\usepackage{amssymb}
\usepackage{amsfonts}  % For additional math symbols
\usepackage{bm}        % Bold math symbols
% Text Formatting and Layout
\usepackage{xcolor}
\usepackage{setspace}    % Allows setting line spacing
\usepackage{float}       % Improved placement of figures
\usepackage{adjustbox}   % Optional for extra control on graphics
\usepackage{makecell}    % Advanced table cell formatting
% Tables and Columns
\usepackage{tabularx,booktabs}
\usepackage{multirow}
\usepackage{multicol}
\usepackage{ulem}
% Listings for Code
\usepackage{listings}
% Algorithms
\usepackage{algorithm}
\usepackage{algorithmic}
% URLs and Hyperlinks
\usepackage{url}         % URL formatting
\usepackage[hidelinks]{hyperref}  % Clickable links in PDF
% Bibliography
\usepackage{natbib}      % BibTeX bibliography style
% Captions and TOC
\usepackage[labelfont=bf]{caption}
\usepackage[titles]{tocloft}
% Additional Packages
\usepackage[section]{placeins}  % Place figures/tables within sections
\usepackage{titlesec}        % Customizable section titles
\usepackage{enumitem}        % Enhanced list formatting
\usepackage{afterpage}       % Code to run after a page break
\usepackage{pdflscape}       % Landscape pages
\usepackage{fancyhdr}        % Fancy headers and footers
%\usepackage{amsaddr}         % AMS address formatting
\usepackage{authblk}
\usepackage{lineno}          % Line numbers
\usepackage{newfloat}        % Define new floating environments
%\usepackage{mdsymbol}      % Conflicts with mathptmx, comment out if needed
% TikZ for Diagrams
\usepackage{tikz}
% CSV Simple for Reading CSV Files
\usepackage{csvsimple}


\newcommand{\fixedfigwidth}{0.95\textwidth}  % Define a consistent figure width

\newcolumntype{P}[1]{>{\centering\arraybackslash}p{#1}}
\newcolumntype{M}[1]{>{\centering\arraybackslash}m{#1}}

% Define a new command \address{}
\newcommand{\address}[1]{%
	\begin{center}
		\large\textit{#1}
	\end{center}
}


\DeclareFloatingEnvironment[name={Supplementary Figure},fileext=lsf,listname={List of Supplementary Figures}]{suppfigure}
\DeclareFloatingEnvironment[name={Supplementary Table},fileext=lsf,listname={List of Supplementary Tables}]{supptable}
\DeclareFloatingEnvironment[name={Supplementary Material},fileext=lsf,listname={List of Supplementary Material}]{suppmat}
\DeclareFloatingEnvironment[name={Supplementary Text},fileext=lsf,listname={List of Supplementary Text}]{supptext}

\linenumbers

\DeclareMathOperator*{\argmax}{arg\,max}
\DeclareMathOperator*{\argmin}{arg\,min}
\newcommand{\indep}{\perp \!\!\! \perp}
\newcommand{\nindep}{\not\!\perp\!\!\!\perp}

\renewcommand\thesection{\arabic{section}}

\begin{document}
	\title{Detecting causal models with Mendelian randomization for genomic trios with many confounding variables}
	\author[1,2,*]{Jarred M. Kvamme}
	\author [3,4,*,$\dagger$]{Audrey Qiuyan Fu}
	
	\affil[1]{Research Computing and Data Services, University of Idaho}
	\affil[2]{Institute for Interdisciplinary Data Sciences, University of Idaho}
	\affil[3]{Department of Family Medicine and Public Health Sciences, Wayne State University School of Medicine}
	\affil[4]{Center for Molecular Medicine and Genetics, Wayne State University}
	\affil[*]{These authors contributed equally to this work.}
	\affil[$\dagger$]{Correspondence: \href{audreyfu@wayne.edu}{\textit{audreyfu@wayne.edu}}}
	\maketitle
	
\begin{doublespace}

%\begin{abstract}
\section*{Abstract}
The principle of Mendelian randomization (MR) makes it possible to infer the potential causal relationship between two biological variables (e.g., expression of two genes, or between two molecular phenotypes such as gene expression and DNA methylation), using a genetic variant as an instrumental variable and thus forming a trio.  Inference is often influenced by many other variables (such as other genes and demographic variables).  Being able to properly account for confounding is therefore essential to causal inference.  However,  methods that can deal with many confounding variables are limited.  Here, we introduce a regression-based method under this principle, namely MR for Genomic Networks in trios (MRGNtrio), to detect diverse causal models for genomic trios using individual-level data while accounting for many confounding variables.  
%Our method detects possible confounders from genomewide data and accounts for them in causal inference.  
We demonstrate the inference power and computational efficiency through extensive simulation and application to examine how a single nucleotide polymorphism potentially regulates its cis- and trans-target genes.
%, and the other is to profile the relationship between DNA methylation and gene expression in breast cancer patients, using copy number alteration as the instrumental variable.  

\setcounter{section}{0}
\section{Introduction}

An important goal in biology is to understand how genes regulate one another and how genetic variation influences gene expression \cite{camacho2018next, lecca2021machine, xu2019machine, pingault2018using}. Gene regulation involves complex networks of high dimensions, which may be represented by directed networks where directed edges represent causal (or regulatory) relationships. Although the ultimate goal in biology is to reconstruct gene regulatory networks at the genome level, it is often of interest to understand the causal relationship between a molecular phenotype (e.g. the expression level) of two genes, or between two molecular phenotypes (e.g., between gene expression and DNA methylation), which may be influenced by many other variables such as other genes, other biological processes, demographic variables, etc.

To learn the (directed) causal relationship between two phenotypes, we can apply the principle of Mendelian randomization (PMR) and leverage the information in genetic variants. The PMR assumes that genotypes in a natural population are randomly assigned, thus observational omics data can be viewed as being collected from a randomized experiment \cite{davey2014mendelian, aten2008using}. This implies that an edge connecting a genotype to a phenotype will always point to the phenotype, but not the other way around \cite{badsha2019learning, yazdani2016generating}. As a result, this constraint on edge direction makes it possible to break the symmetry of association and draw inferences about the regulatory structure among phenotypes \cite{badsha2019learning, ludl2021comparison, davey2014mendelian}.   

Many methods have been developed under the PMR \cite{wang2014.network, wang2017efficient, millstein2009cit, zeng2021statistical, badsha2019learning}, although they often are interested in the mediation model $V \rightarrow T_1 \rightarrow T_2$, where $V$ is the genetic variant, $T_1$ and $T_2$ are the variables of interest. In this model, the mediator $T_1$ relays the regulatory information from $V$ to the outcome $T_2$; for example, a genetic variant regulates the expression level of a transcription factor, and its reduced expression leads to its downstream target gene being under-expressed. However, other causal relationships may exist among the three variables. 

Previously \cite{badsha2019learning} developed MRPC, a method that goes beyond mediation and infers five basic causal models for a trio that are mutually exclusive (\autoref{fig:trio}) \cite{badsha2019learning, badsha2021mrpc}. They include a null model ($V \rightarrow T_1$ and no relationship between $T_1$ and $T_2$, denoted M0) and four non-null models: mediation (denoted M1), v-structure (or reverse causation; $V \rightarrow T_1 \leftarrow T_2$, denoted M2), conditional independence (or horizontal pleiotropy; $T_1 \leftarrow V \rightarrow T_2$, denoted M3), and a fully connected model ($T_1 \leftarrow V \rightarrow T_2$ and $T_1 \leftrightarrow T_2$, denoted M4). 

The true model behind M4 may be directed: the edge between $T_1$ and $T_2$ is $T_1 \rightarrow T_2$ or $T_1 \leftarrow T_2$.  However, these two directions have the same likelihood and are therefore Markov equivalent. Without additional information, they are indistinguishable and only an undirected edge can be inferred.   

The PMR and the five basic causal models require that no confounding variables can explain away the directed edge between $T_1$ and $T_2$, since this edge is our main interest.  However, as mentioned earlier, many variables can influence the inference due to the complexity of gene regulation.  A confounding variable may influence the inference in different ways.  \cite{yang2017identifying} proposes three types of confounding for the $T_1-T_2$ relationship: as a confounder, an intermediate variable, or a common child to $T_1$ and $T_2$. Specifically, a confounder is a common parent to both $T_1$ and $T_2$ and can induce a false edge between $T_1$ and $T_2$ when unaccounted for.  This is similar to $T_1$ and $T_2$ being marginally dependent in M3.  

An intermediate variable acts like a mediator between $T_1$ and $T_2$ and can potentially explain away a true $T_1 \rightarrow T_2$ edge when this intermediate variable is accounted for in the inference, leading to a false negative conclusion.  By contrast, a common child can induce a false edge between $T_1$ and $T_2$ when accounted for, leading to a false positive conclusion.  This is because when the common child is accounted for in the inference, $T_1$ and $T_2$ become conditionally dependent, similar to their relationship in M2.  

Accounting for confounding variables is therefore crucial in causal network inference \cite{burgess2015network, yang2017identifying, shan2019identification, amar2021graphical, zhao2024adjusting}. Two basic approaches have emerged to address latent confounding between an exposure and outcome. In the first approach, candidate confounding variables are incorporated into the graph as additional nodes in the network and also treated in the same way as other nodes; our MRPC method and several other methods take this approach \cite{badsha2019learning, amar2021graphical}. However, it is computationally intensive to include many confounding variables, as the computational complexity for graph inference grows super exponentially with the size of the graph. The second approach focuses on identifying and accounting for confounders \cite{yang2017identifying, shan2019identification, burgess2015network}. This approach includes all the confounders as a whole in the conditional independence tests, and does not increase the size of the network.  For example, we can include the confounders when testing the $T_1 - T_2$ edge in a trio.  The trio remains a network of three nodes, but the tests are conditioned on these confounders.  However, these methods focus on the mediation model.

Herein we present Mendelian Randomization for Genomic Networks in trios (MRGNtrio), 
a novel machine learning method for inferring the causal model of a trio from individual-level data in the presence of potentially hundreds of confounding variables.  MRGNtrio builds on MRPC and infers the five basic causal models.   We take the second approach to confounding described above and focus on accounting for confounders, and not other types of confounding variables.  This approach allows for the inclusion of many confounding variables all at once, while maintaining high power and fast computation times. We demonstrate the performance of MRGNtrio and compare to other methods through extensive simulation studies.  We further assess this method in our study of the cis- and trans-regulatory modes of genetic variants, using the whole blood data from the Genotype-Tissue Expression (GTEx) Consortium \cite{gtex2020}. 

\section{Methods}
\subsection{The MRGNtrio model}

The five causal models for a trio can each be described by a set of marginal and conditional dependence relations (\autoref{tab:tests}). Specifically, in the mediation model (M1), any two variables are dependent on each other, but $T_2$ and $V$ are conditionally independent given $T_1$.  In the v-structure (M2), $T_2$ and $V$ are marginally independent but conditionally dependent given $T_1$. This is because when the value of $T_1$ and that of $T_2$ (or $V$) are known, the possible values of $T_1$'s other parent are constrained.  

In the conditional independence model (M3), $T_1$ and $T_2$ are marginally dependent but conditionally independent given $V$. In the fully connected model (M4), all pairs of variables are marginally and conditionally dependent. By contrast, in the null model (M0), if $V$ is connected to $T_1$ (or $T_2$), then the third variable $T_2$ (or $T_1$) is marginally and conditionally independent of $V$ and of $T_1$ (or $T_2$). These relations indicate that M0, M1 and M3 can be uniquely determined by conditional (in) dependencies, and that M2 and M4 can be distinguishable from M0, M1 and M3, although marginal dependences are needed to further separate M2 and M4.

MRGNtrio leverages this observation and prioritize statistical tests for conditional dependences, which also allows us to account for a set of confounders, denoted by the vector ${\bf U}$, through the following regression: 
%
\begin{align}
	T_1 &= \alpha_1 +\beta_{11}V+\beta_{12}T_2+{{\bf \Gamma}_1 {\bf U}}+ \epsilon_1; \label{eq:reg1} \\
	T_2 &= \alpha_2 +\beta_{21}V+\beta_{22}T_1+{{\bf \Gamma}_2 {\bf U}}+ \epsilon_2, \label{eq:reg2}
\end{align}

\hspace{-17pt}where $\alpha_i$s are the overall means, scalar $\beta_{ij}$s and vectors ${\bf \Gamma}_i$s are effect sizes, and $\epsilon_i$s are random errors that each follow a normal distribution with mean 0.  The conditional (in) dependencies are encoded in the $\beta$s.  For example, the null hypothesis $H_0: \beta_{11} = 0$ 
represents conditional independence between $T_1$ and $V$ given $T_2$ and confounding variables ${\bf U}$, which is denoted as $T_1 \indep V | \{T_2, {\bf U} \}$, whereas the $H_0: \beta_{21} = 0$ 
represents conditional independence $T_2 \indep V | \{T_1, {\bf U} \}$.  If we do not reject $H_0: \beta_{21} = 0$, which means that $T_2$ is conditionally independent of $V$, then both M1 and M0 are possible models.  If we further reject $H_0: \beta_{22} = 0$, which means that there is still association between $T_2$ and $T_1$, then we infer this trio to follow an M1 model.

M2 and M4 are further distinguished by the marginal relationships between the genetic variant and each phenotype. In the v-structure M2, if $V$ is connected to $T_1$ (or $T_2$), then $V$ and $T_2$ (or $T_1$) are marginally independent.  In M4, however, $V$ has marginal dependence on both $T_1$ and $T_2$.  The marginal relationship can also be encoded in the coefficient $\beta$ in the following univariate regression:
%
\begin{align}
	T_i &= \alpha + \beta_{i1} V + \epsilon, \;\;\;\; \text{$i = 1$ or 2}. \label{eq:eq3}
\end{align}

\hspace{-9pt} Marginal independence is represented by $\beta=0$. When the data is not informative enough for any of the basic models to be inferred, we categorize the trio as ``other".

We perform $t$ tests for the hypotheses on the regression coefficients. We also allow for a permutation test for coefficients $\beta_{12}$ and $\beta_{22}$ when the normality assumption for the regression and $t$ test is questionable, because our primary interest is the relationship between $T_1$ and $T_2$ and also because the relationship between $V$ and either $T$ node is easier to infer.  Permutation is done within each level (i.e., genotype here) of the categorical variable $V$, similar to \cite{yang2017identifying}. 


For example, to test $\beta_{12}$ in regression (\autoref{eq:reg1}), we permute the values of $T_1$ within each genotype and calculate the permutation $t$ statistic $t_{12}^*$; this procedure removes any dependence between $T_1$ and $T_2$ while preserving the possible dependence between $T_1$ and $V$.  The collection of the permutation $t$ statistics generates a null distribution for $t_{12}$.  We compute a two-sided permutation $p$-value as 
\begin{align}
	p_{12} = 2P\left( Z > \frac{|t_{12}^{obs}|-\mathbb{E}\left[t_{12}^*\right]}{\sqrt{\text{Var}\left(t_{12}^*\right)}}\right)\label{eq:perm_p} ,
\end{align}

\hspace{-9pt}where $\mathbb{E}$ denotes the expected value of $t_{12}^*$ under the null distribution, and $Z$ is the standard normal random variable.


\subsection{Confounder selection}
As part of the MRGNtrio model, we adopt a similar adaptive confounder selection strategy to that described in \cite{yang2017identifying}. Given a pool of potential confounders, we first exclude variables that are marginally associated with $V$. This filtering step is motivated by the assumption that true confounders are parents of the $T$ nodes and are therefore independent of each other and of $V$. In contrast, variables downstream of one or both $T$ nodes—such as common children or intermediates—may appear associated with $V$ due to collider bias or mediation effects, and should not be treated as confounders.

Next, candidate confounders are selected based on their statistical association with either $T_1$ or $T_2$. Following the approach of \cite{yang2017identifying}, we test whether the regression coefficients from modeling each candidate confounder on both $T_1$ and $T_2$ are simultaneously zero. This step involves multiple hypothesis testing, and both the initial filtering and confounder selection stages employ the $q$-value procedure to control the false discovery rate (FDR) \cite{storey2003positive}.

While we retain the same statistical framework as \cite{yang2017identifying}, we substantially improve computational efficiency in two key ways. First, we parallelized the confounder filtering and regression-based selection hypothesis tests across candidate confounders, allowing for faster evaluation when the pool of confounders is large. Second, we leverage the R package \textit{propagate} to improve correlation matrix computations by partitioning the problem into smaller, memory-efficient blocks \cite{spiess2018package}, which scales more efficiently to large genomic datasets.


\subsection{Simulation study}
We examined the performance of MRGNtrio in a wide range of settings and in comparison to similar methods; MRPC \cite{badsha2019learning} and GMAC \cite{yang2017identifying}. We simulated $1,500$ trios by sampling from the five causal models with equal probability. This resulted in approximately $300$ trios of each model. For each trio, we applied different combinations of effect size between nodes, noise, and frequency of the minor allele of the genetic variant. We fixed the sample size for all trios at $n = 670$, which is the same sample size of the whole blood tissue in the Genotype-Tissue Expression (GTEx) Consortium (version 8) to which we apply our method \cite{gtex2020}.  We also included $1$-$15$ confounders (median: 8), one common child and one intermediate variable in each trio so that we can assess confounder selection and the impact of confounder misspecification on causal inference.      

When analyzing these trios, we assumed that all the confounding variables were available but we did not know which confounding variables were associated with which trio. In our inference, we first performed the confounding variable selection method of \cite{yang2017identifying} to all the confounding variables, before applying three methods: MRGNtrio, GMAC and MRPC to each trio with the identified confounders. For all these methods, we used the $q$-value \cite{storey2003positive,storey2002qvalue} method for FDR control in confounder selection (denoted CSFDR). For MRGNtrio, we also performed confounder selection without the FDR control (denoted CSnoFDR) (see ``Application of causal inference methods").We further simulated another set of 300 trios with 15–50 confounders (median: 34), representing a scenario with moderate to to extreme genetic confounding. The candidate confounder pool trios in this scenario contained about $11,000$ variables.

We simulate genetic trios consisting of a bi-allelic genetic loci with three representative genotypes and two expression target genes under the five causal structures (\autoref{eq:snpsim}). Specifically, for models M0 - M2 which each have two possible configurations, we simulate M0 as $V \rightarrow T_1$; no relationship between $T_1$ and $T_2$, M1 as $V\rightarrow T_1 \rightarrow T_2$, and M2 as $V\rightarrow T_1 \leftarrow T_2$. Furthermore, since the undirected edge between $T_1$ and $T_2$ in model M4 can be represented by either $T_1 \leftarrow V \rightarrow T_2; T_1 \rightarrow T_2$ or $T_1 \leftarrow V \rightarrow T_2; T_2 \rightarrow T_1$ we restrict all simulated M4 trios to the former configuration for simplicity. 

For all trios, we vary the minor allele frequency, signal strengths of each effect represented by the coefficients in (5), and the dispersion of the residuals represented by $\sigma$. Specifically, we fix the effect from the genetic variant to be the same for all edges with the genetic variant (i.e $V \rightarrow T_1$ or $V \rightarrow T_2$ ), and sample the strength of genetic effect uniformly between $0.5$ and $1.5$ \cite{yang2017identifying}. 

Similarly, we fix the effect between the genes and sample the strength of the effect between the genes (i.e the $T_1 - T_2$ edge) uniformly between $0.5$ and $1$. Lastly, We simulate the strength of the residual standard deviation $\sigma$ by scaling the effect of the $T_1 - T_2$ edge such that $\sigma = c\beta$. We sample the scaling effect $c$ uniformly between $0.3$ and $1.5$

We utilize the topological ordering of the graph of each trio to determine the order in which nodes are generated. Under this framework, the eQTL is simulated first from a multinomial distribution under Hardy-Weinberg equilibrium:
%
\begin{align}
	Pr(V_i = 0 ) = (1-\theta)^2; \ Pr( V_i = 1) =2\theta(1-\theta); \ Pr(V_i = 2) = \theta^2, \label{eq:snpsim}
\end{align}

\hspace{-5pt}where $\theta$ represents the frequency of the minor allele of the genetic variant and $V_i$ is the genotype of the $i^{th}$ individual. We simulate the minor allele frequency uniformly between $0.01$ and $0.5$ so that approximately $10\%$ of simulated trios have a ``rare" allele.  

For the remaining nodes, we simulate each gene or confounding variable $Y_j$ as coming from a normal distribution centered at the linear combination of its parent nodes $S = {\bf Parents}(Y_j)$ and with the intercept at the origin (\autoref{eq:sim_reg}). The set of parents $S$ may include the genetic variant, the other gene, confounding variables, a combination, or it may be the empty set. 
\begin{align}
	Y_j &\sim N\left( \beta_0 + \sum_{i\in S} \beta_{i} Y_i, \ \sigma^2 \right) \label{eq:sim_reg}
\end{align} 

We consider three types of latent confounding variables that may be present for a trio; confounders, intermediate variables, and common child variables \autoref{fig:confounding}) \cite{yang2017identifying}. We consider only confounding variables that impact the association between target traits of the genetic variant ($T_1$ and $T_2$), as it has been previously established that there are likely few variables that confound the relationship between a genetic variant and its target molecular traits \cite{davey2014mendelian, burgess2018inferring}. 


All trios in our simulation are simulated to have a number of latent confounding variables that must be identified through statistical inference. Specifically, each trio is simulated to have a varying number of confounders, one common child variable, and one intermediate variable. All confounding variable effects are simulated independently from a uniform distribution. 
Following \cite{yang2017identifying}, we set confounder and intermediate variable effects to range between 0.15 and 0.5 and common child effects between 0.5 and 1.5. Additionally, we allow all confounding variable effects a $50\%$ probability for the effect to be negative. 

\subsubsection{Performance metrics}

We summarized the inference results by averaging the performance metrics across all simulated trios using three sets of measures (see ``Assessing performance on simulated trios"). First, we calculated recall and precision to assess whether the causal model type (i.e M0, M1, etc.) is correctly inferred for each trio computed as:
\begin{align}
	\text{Recall}_\text{model} &= \frac{\text{Number of correctly inferred trios}}{\text{Total number of trios of the causal model of interest}}; \label{eq:rec_total}\\
	\nonumber \\
	\text{Precision}_\text{model} &= \frac{\text{Number of correctly inferred trios}}{\text{Total number of trios inferred to be of the causal model of interest}}. \label{eq:prec_total}
\end{align}


\noindent This performance measure tends to be conservative as a good performance requires the entire trio to be inferred correctly. 

Second, we calculated recall and precision based on individual edges in a trio.  Similar to our earlier work \cite{badsha2019learning}, we assigned a score of 0.5 if a true edge is inferred to be present, and another score of 0.5 if the inferred edge has the correct direction. We summarize the edge-based performance across all trios by averaging the individual precision and recall scores: 

\begin{align}
	\text{Recall}_\text{edges} &= \frac{\text{Number of correct edges in inferred graph}}{\text{Total edges in true graph}}; \label{eq:rec_edges}\\
	\text{Precision}_\text{edges} &= \frac{\text{Number of correct edges In inferred graph}}{\text{Total edges in inferred graph}}.\label{eq:prec_edges}
\end{align}

Lastly, we calculated recall and precision based on individual edges in a trio (\autoref{eq:rec_edges}-\ref{eq:prec_edges}).  We assigned a score of 0.5 if a true edge is inferred to be present, and another score of 0.5 if the inferred edge has the correct direction.	We summarize the edge-based performance across all trios by averaging the individual precision and recall scores.

Lastly, we incorporate a performance measure strictly for the presence/absence of the edge between the genes ($T_1 - T_2$ edge) to allow for direct comparison of our method with GMAC:
\begin{align}
	\text{Recall}_\text{$T_1 - T_2$} &= \frac{\text{Number of trios with a correctly identified $T_1 - T_2$ edge}}{\text{Total trios with a true $T_1 - T_2$ edge}}; \label{eq:rec_tedge}\\
	\nonumber \\ 
	\text{Precision}_\text{$T_1 - T_2$} &= \frac{\text{Number of trios with a correctly identified $T_1 - T_2$ edge}}{\text{Total trios with an inferred $T_1 - T_2$ edge}}. \label{eq:prec_tedge}
\end{align}


\noindent This metric confines the five causal models to two classes: models with $T_1 - T_2$ edge present (M1, M2, M4) and models with $T_1 - T_2$ edge absent (M0, M3). 

\subsection{Application of causal inference methods}

We compare our MRGNtrio method to two related PMR-based causal inference methods: Genomic Mediation analysis with Adaptive Confounding adjustment (GMAC) \cite{yang2017identifying} and Mendelian Randomization PC (MRPC) \cite{badsha2019learning}. Unlike MRGNtrio and MRPC, GMAC is a medation - based method that is restricted to inferring only the causal edge between $T_1$ and $T_2$. To that end, we only compare MRGNtrio and MRPC to GMAC on the performance for predicting the $T_1 - T_2$ edge. We set the number of permutations to $1000$ and the test threshold for all marginal and conditional tests to be $\alpha = 0.01$ which is approximately a Bonferroni correction. For MRPC, we use the ADDIS FDR control, Gaussian CI test, and set the FDR to $5\%$. For MRGNtrio, we employ the permutation test whenever the minor allele is observed in less than $5\%$ of samples. GMAC, however, applies the permutation test described in to all trios regardless of allele frequency. For both methods, we set the number of permutations to $1000$. 

Following the approach of \cite{kvamme2022types}, we run GMAC twice on all trios, first with $T_1$ as the mediator and then with $T_2$ as the ``mediator" and take significant mediation at the $\alpha < 0.05$ and $\alpha < 0.01$ levels. To give GMAC the best chance at predicting the $T_1 - T_2$ edge across all causal models, we combine the inference results from the tests with $T_1$ as the mediator with the tests with $T_2$ as a mediator. Combining the inference from using each gene as the mediator establishes four possible results: positive with $T_1$ as the mediator, positive with $T_2$ as the mediator, both positive, or both negative. To afford GMAC greater power across the five models, we take any of the aforementioned positive cases for mediation as an indication of the existence of the $T_1 - T_2$ edge. 

To emulate a real statistical analysis, we include confounding variable selection as a preliminary step to inference. The confounder selection method of \cite{yang2017identifying} requires the input of a pool of potential confounding variables. In simulation, we construct this pool of potential confounding variables by aggregating the true confounding variables from all simulated trios together (i.e all confounders, intermediate and common child variables). We control the false discovery rate (FDR) for filtering common child and intermediate variables from the candidate pool at the $10\%$ level. Furthermore we consider two settings for selecting confounders from the pool post-filtering (i) selecting confounders with no FDR correction and significant confounders selected at the $\alpha < 0.01$ (denoted CSnoFDR) and (ii) using the qvalue method to control the FDR of selected confounders at the $5\%$ level (denoted CSplusFDR). The latter setting is the default for selecting confounders using the GMAC software. 


\subsection{Application to GTEx data}

The GTEx consortium has mapped a large number of eQTLs and their associated cis target genes (cis-eQTLs and cis-target genes) for most genes in the human genome across many human tissue types \cite{gtex2020}. Recent  studies have focused on identifying cis-target genes that mediate the relationship between the eQTL and its associated trans-target genes \cite{yang2017identifying}.  We previously have applied the MRPC method to all the GTEx tissues to detect diverse causal models for such trios: an eQTL and its associated cis- and trans-target genes, where the genes are protein coding genes or long non-coding RNAs (lncRNAs) \cite{kvamme2022types}.  

We implement the preprocessing and normalization pipeline from Section 2.2 to construct trios for WholeBlood tissue from version 8 of GTEx. Each trio consists of an eQTL and its highest association ``cis" and ``trans" target genes.  We define cis genes as genes less than one megabase from the eQTL and trans genes as genes residing more than one megabase from the eQTL or on a different chromosome. As noted in Section 2.2 individual trios may share the same eQTL or the same gene, and the gene may be a cis-gene in one trio but a trans gene in another trio. 

In GTEx version 8, the number of available samples for whole blood is 670. After preprocessing, 8,823 trios were identified. To eliminate bias induced by mismapping transcripts, we restricted our final set of trios to protein coding genes or long non-coding RNA genes \cite{yang2017identifying}. This additional filtering step reduces the number of trios to 3,248.  Many of eQTLs provided by the GTEx consortium contain missing values. Since the GMAC algorithm is unable to handle missing values, we used multiple correspondence analysis implemented in the \textit{missMDA} package in $R$ \cite{josse2016missmda} to impute the missing genotypes of individuals using the collection of unique eQTLs from all GTEx trios analyzed for whole blood. 

\subsubsection{Confounding variable selection and filtering for GTEx trios}

We performed principal component analysis (PCA) on the PEER-normalized, genome-wide gene expression data and used the resulting principal component (PC) scores as the pool of candidate covariates for confounder selection. Each PC summarizes the regulatory effects of a large number of genes across the genome. Because each trio represents a small component of a larger regulatory network, PCs that correlate with a trio may capture potential confounding effects arising from pleiotropy. We performed confounder selection as described in Confounding Variable Selection and applied common-child and intermediate-variable filtering to obtain a final set of confounders for each trio. Following our findings from simulation, we applied MRGNTrio using the CSnoFDR setting and selected confounders at a significance threshold of $\alpha < 0.05$. While the simulation study used a more stringent cutoff of $\alpha < 0.01$, we relaxed this threshold here because the more stringent cutoff produced nearly the same set of confounders as the CSFDR setting when applied to the whole-blood PC score matrix. For each trio, we also included sex, platform, and PCR batch as additional covariates.

\section{Results}

\subsection{Simulation Results}
\subsubsection*{MRGNtrio improves power and precision in the presence of confounding}

MRGNtrio achieves substantially higher recall for the mediation, v-structure, and fully connected models than MRPC does (M1: $+82\%$, M2: $+47\%$,  and M4: $+25\%$; (\autoref{tab:results}). All these models contain a $T_1 - T_2$ edge, indicating the sensitivity of MRGNtrio in detecting the $T_1 - T_2$ edge. On the other hand, this sensitivity leads to reduced power to detect the null and conditional independence models (M0: $-37\%$, M3: $-38\%$), where there is no edge between the $T$ nodes. Remarkably, MRGNtrio also achieves equivalent or better precision for nearly all models (M0: $+47\%$, M1: $+29\%$, M2: $+48\%$, M3: $+8\%$, and M4: $-17\%$. 

When comparing the performance using edge-level recall and precision, MRGNtrio infers more edges than MRPC while maintaining a high degree of precision. Furthermore, MRGNtrio matches the performance of GMAC for the $T_1 - T_2$ edge (\autoref{fig:main_t1_t2_inf}), as the GMAC test for mediation is a special case of our method. The similar performance also shows that the notion of ``mediation" detected by GMAC is broad, while our method is able to distinguish the different models behind the broadly-defined mediation. 

\subsubsection*{MRGNtrio handles many confounding variables}

In scenarios with moderate to severe confounding, MRGNTrio and GMAC performed similarly well in recall and maintained consistently high recall for the $T_1 - T_2$ edge, even under substantial confounding, with precision remaining relatively unchanged (\autoref{fig:main_t1_t2_inf}). Additionally, recall and precision across all edges remained high, indicating that MRGNTrio can accurately infer most trio edges, except in cases with many ($>35$) confounders, where recall declined but still exceeded that of MRPC \autoref{fig:main_all_edge}.

\subsubsection*{The impact of confounder selection on causal network inference}

Accurate selection of confounding variables is a critical step preceding causal inference in Mendelian randomization studies, as errors at this stage can propagate to the causal model by omitting informative genetic or environmental covariates. In our simulation study, we compared two strategies for selecting confounders from a pool of candidate covariates: CSnoFDR, which selects confounders without false discovery rate (FDR) correction, and CSFDR, which controls the FDR at 5\%. We first evaluated both methods based on their ability to identify the true confounders, measured by precision and recall. Overall, both methods achieved recall greater than 90\% when the number of true confounders per trio was modest (e.g., fewer than 25; see \autoref{fig:SF4_Confounder_Selection_Performance}). However, CSFDR showed reduced performance when for trios with many true confounders, with average recall to detect true confounders dropping to approximately 65\%. Conversely, CSnoFDR achieved nearly perfect recall across all simulations, but at the expense of very low precision, indicating that it selected many false confounders in addition to the true ones.

When we examined how confounder selection affected causal inference performance, we found no clear relationship between the number of false-positive confounders and MRGNTrio’s ability to accurately infer causal edges (\autoref{fig:SF7_Confounder_Selection_Impact_On_Inference}). In contrast, falsely omitted confounders had a strong negative impact on causal inference performance. Overall, we observed little difference in causal inference performance between the two methods in our simulation study. This may be because only a small number of trios were simulated with a large number of confounding variables (e.g., more than 35), limiting our ability to detect any differences between the two methods. Although CSnoFDR is more liberal, it may be preferable in scenarios where many potential covariates are expected for a trio.


\subsubsection{Speed comparison of MRGNtrio, MRPC and GMAC}

Our analysis involves two major steps: confounder selection and causal network inference. We report the computation time for each step separately. We conducted a comparative analysis of the time required for confounder selection across varying sizes of the confounder covariate pool (100, 500, or 1000 covariates) and different numbers of trios for which confounders needed to be identified (500 or 1000), resulting in a total of six replicates. Notably, in the most intensive case 1000 covariates and 1000 trios, our method accomplished confounder selection within 16 minutes (\autoref{fig:SF1_comp_times_for_conf_selection_updated}). This represents a notable gain in computation efficiency compared to the implementation provided by GMAC, which took approximately 28 minutes. Although our confounder selection procedure is similar to that of \cite{yang2017identifying}, our implementation is more efficient due to the use of the R package propagate \cite{spiess2018package}, which enables fast correlation calculations between each trio’s target genes and the candidate covariate pool, significantly reducing the time required for confounder selection.

In causal network inference, MRGNtrio is substantially faster than both MRPC and GMAC (\autoref{tab:time_comp}). The median time to analyze a trio using the same set of confounders was 0.0065 seconds for MRGNtrio, compared to 0.3557 seconds for MRPC and 0.7088 seconds for GMAC. This makes MRGNtrio over 50 times faster than MRPC and more than 100 times faster than GMAC, on average.

The extended inference time for MRPC is largely due to the ADDIS method used for false discovery rate (FDR) control, which was particularly slow for some trios. When MRPC is applied using the more conservative LOND method for FDR control, its runtime becomes comparable to that of MRGNtrio. For GMAC, we note that our comparison used the default single-threaded implementation provided in its R package. Enabling the parallelization option available in GMAC can substantially reduce its computational time.








\subsubsection{Confounding contributes to an inflated type I error rate}

In our simulation study, we observed significantly higher rates of false positive mediation for the GMAC method than those reported by \cite{yang2017identifying}. In their simulations under the null model, \cite{yang2017identifying} reported type I error rates (false mediation) consistent with the significance level used to test for mediation. However, in our simulations, we observed type I error rates for the $T_1 - T_2$ edge under the null model to be as high as $75\%$ for both MRGNtrio and GMAC at the $\alpha < 0.01$ significance cutoff. We attribute these differences to several distinctions between our simulation framework and theirs. 

First, we simulated a substantially larger number of confounders for each trio (1--50 across all trios in our simulations, compared to 1--2 in \cite{yang2017identifying}). Additionally, to construct the pool of confounders used for selection, we aggregated all confounding variables from all simulated trios, resulting in a large candidate pool of more than $10{,}000$ variables. We also simulated each trio to have one common child and one intermediate variable, which were included in the candidate covariate pool. In contrast, \cite{yang2017identifying} simulated a much smaller pool of $350$ confounders, yielding an average of approximately two confounders per trio, with many trios having no confounders. Furthermore, their simulations treated each type of variable (confounder, common child, intermediate variable) separately, whereas our simulations included all three types simultaneously.



%To validate the error rate reported \cite{yang2017identifying} we integrated their simulation with ours, examining the effect of confounders on Type I error (false mediation) under the null model. Specifically, we simulated 100 additional trios under the null model for each of seven scenarios: 2, 4, 6, 8, 10, 12 or 15 confounders, excluding common child and intermediate variables. We then applied the GMAC method with CSFDR confounder selection to all trios in each scenario but using our approach of aggregating covariates across simulated trios to form confounding variable pool.  The first scenario with 2 confounders is comparable to the simulations under the null model presented by \cite{yang2017identifying}, while the remaining scenarios are similar to our simulations. 
%
%Our results show that we are able to replicate the type I error rates reported by \cite{yang2017identifying} in the first scenario with 2 confounders (\autoref{fig:SF9_GMAC_Validation_Summary}). All other scenarios show increasingly inflated rates of false mediation that scale proportionally to the number of confounders and the size of the covariate pool. Thus significant confounding leads to significantly inflated rates of false positive mediation and further supports the low power to detect the null and conditional independence models observed in our simulation studies.

To validate the error rates reported by \cite{yang2017identifying}, we incorporated their simulation design into our framework to examine the effect of confounders on type~I error (false mediation) under the null model. We simulated 100 additional trios for each of seven scenarios containing 2, 4, 6, 8, 10, 12, or 15 confounders, excluding common child and intermediate variables, and applied GMAC with CSFDR confounder selection using our approach of aggregating covariates across simulated trios to construct the confounding variable pool. The scenario with two confounders is comparable to the null-model simulations reported by \cite{yang2017identifying}, while the remaining scenarios align with our simulation settings. We successfully replicated the type~I error rates reported by \cite{yang2017identifying} in the two-confounder scenario (\autoref{fig:SF9_GMAC_Validation_Summary}), whereas all other scenarios showed increasingly inflated rates of false mediation that scaled with the number of confounders and the size of the covariate pool. These results demonstrate that substantial confounding leads to markedly inflated false positive mediation and further support the low power to detect the null and conditional independence models observed in our simulation study.




\subsection{Application results}

We applied MRGNTrio combined with the CSnoFDR confounder selection procedure to $3{,}248$ eQTL trios for whole blood in GTEx. We compared the results from MRGNTrio with those from MRPC (using the ADDIS method to control the overall FDR of conditional tests) and GMAC \cite{badsha2019learning,yang2017gmac}. Principal component analysis (PCA) was performed across all genes in whole blood to derive $670$ principal components, which served as the pool of candidate confounding variables. For each trio, we identified a median of $51$ confounders (SD = $6.38$) (\autoref{fig:main_GTEx}). 

After applying MRGNTrio, $911$ trios were classified as the null model, $91$ as mediation, $10$ as v-structures, $2,062$ as conditionally independent, and 104 as fully connected. The remaining $70$ trios were categorized as “other.” We compared these results with our previous work using MRPC \cite{kvamme2022types}, and found that MRGNTrio identified more than twice as many cases of causal mediation and $173$ additional cases of conditional independence. Similar to MRPC, MRGNTrio also distinguishes two types of mediation: (1) cis-gene mediation $V \rightarrow T_1 \rightarrow T_2$, where the cis-gene mediates the signal from genetic variant to the trans-gene, and (2) trans-gene mediation $V \rightarrow T_2 \rightarrow T_1$, where the trans-gene acts as the mediator between the genetic variant and the cis-gene. Among the $91$ trios identified as mediation by MRGNTrio, $69$ (75.5\%) had the cis-gene  ast the mediator and $22$ (24.5\%) had the trans-gene as the mediator (\autoref{tab:S7_GTEx_comp_MRGN_MRPC_GMAC}). Compared to MRPC, MRGNTrio identified not only more mediation trios overall, but also a higher proportion of cis-gene mediation. However, only $8$ cis-gene mediation trios and $6$ trans-gene mediation trios were identified by both methods.

We further investigated trios where MRGNTrio inferred mediation but MRPC did not and found that these discrepancies were primarily driven by differences in the sets of selected confounders (see \autoref{txt:supptext_document}; Examples 1–2). In two illustrative cases, MRGNTrio selected 40 and 43 confounders, whereas MRPC selected only three and none. When we reran MRGNTrio using the same confounder set and inference model as MRPC, both methods produced identical results, confirming that accurate identification of confounders is essential for consistent causal inference.


\subsubsection{Permutation testing causes conservative model shifts among GTEx trios}

Comparing MRGNTrio and GMAC results on GTEx trios, we observed minimal overlap in the identified cases of mediation, similar to our comparison with MRPC. Of the 57 trios identified as mediation by GMAC, only four (two cis-gene mediated and two trans-gene mediated) were also inferred as mediation by MRGNTrio. Two key factors likely contribute to this difference: (i) GMAC uses the CSFDR method for confounder selection, whereas MRGNTrio was applied using CSnoFDR, and (ii) GMAC applies a permutation test to all trios, whereas MRGNTrio applies it only to trios with rare alleles.

Regarding confounder selection, GMAC adjusted for a median of 30 confounders, whereas the more liberal CSnoFDR method used for MRGNTrio adjusted for a median of 51. This highlights a substantial difference in the inference models applied by the two methods.

To assess the effect of the permutation test, we applied MRGNTrio using permutation-based conditional independence tests across all GTEx trios. This changed the inferred model for 89 of the 3,248 trios. Trios affected by permutation had significantly lower minor allele frequencies ($\leq 14\%$) compared to those unaffected ($t=-7.69$, $p=7.2\times10^{-13}$; \autoref{fig:SF3_GTEx_Permutation_and_Model_Changes}). Approximately 97\% of permutation-induced changes were conservative, typically involving edge removal or reclassification as ``Other'' (see \autoref{txt:supptext_document}; Examples 3–4). Similar conservative behavior was observed in our simulation study, although no significant correlation with allele frequency was detected (\autoref{fig:SF5_Simulation_Permutation_and_Model_Changes}).

Finally, we examined the impact of allele frequency on gene expression in GTEx trios. Low-frequency alleles were generally associated with right-skewed trans-gene expression, while cis-gene expression remained approximately normal (\autoref{fig:SF8_Trio_Skew_Stat_Distribtuions}). Trios affected by permutation exhibited higher skewness and kurtosis, reflecting residual non-normality despite normalization efforts. This non-normality was typically driven by a small number of individuals with extreme trans-gene expression, as illustrated in \autoref{fig:trio_number_730}. In these cases, the permutation test mitigated the influence of such outliers, suggesting it serves as an important conservative measure for trios with rare alleles.


%Comparing MRGNtrio results with those of GMAC on GTEx trios, we found minimal overlap in the identified cases of mediation, similar to our comparison with MRPC. Out of the $57$ trios identified as mediation by GMAC, only four (two cis-gene mediated and two trans-gene mediated) were also inferred as mediation using MRGNtrio. It's important to highlight two key differences between our analysis using MRGNtrio and the analysis with GMAC: (i) GMAC utilizes the CSFDR setting for confounder selection, and (ii) GMAC uses the permutation test described at the end of the section ``The MRGNtrio model" to test mediation while MRGNtrio only applies the permutation test to trios with rare alleles. 
%
%To better understand the sources of discrepancy between the GMAC and MRGNtrio results, we examined differences in their confounder selection procedures and the use of permutation testing. GMAC, which employs the CSFDR method, adjusted for a median of 30 confounders, whereas our analysis with MRGNtrio, using the CSnoFDR method, adjusted for a median of $51$. To assess the impact of the permutation test, we applied MRGNtrio using permutation-based conditional independence tests across all trios, which altered the inferred model for 89 of the 3,248 trios. Trios whose inferred models changed under the permutation test had significantly lower minor allele frequencies ($\leq$ 14\%) compared to those whose models remained unchanged ($t = -7.69$, $p = 7.2 \times 10^{-13}$, $df = 191$; see \autoref{fig:SF3_GTEx_Permutation_and_Model_Changes}). Furthermore, approximately 97\% of the model changes induced by permutation testing were conservative, typically involving edge removal or reclassification into the “Other” category (see Supplementary Text; Examples 3-4). Similar conservative behavior of the permutation test was also observed in simulation although it was noted among a much smaller number of trios (see \autoref{fig:SF5_Simulation_Permutation_and_Model_Changes}). 
%
%Further examination of the impact of allele frequency on cis- and trans-gene expression indicated that low-frequency alleles were associated with right-skewed trans-gene expression, while cis-gene expression remained approximately normal (\autoref{fig:SF8_Trio_Skew_Stat_Distribtuions}). All of the trios affected by permutation exhibited higher skewness and kurtosis, suggesting these trios exhibited residual non-normality despite efforts our to normalize it. This non-normality was generally characterized by a small number of individuals with extremely different trans gene expression. We provide an example of one such case in (\autoref{fig:trio_number_730}). In these cases, the permutation test mitigated the influence of such outliers, indicating it may be an important conservative measure in the presence of rare or infrequent alleles.

\subsubsection{MRGNtrio is robust to model misspecification in causal mediation}

Based on our simulation and application findings, we conducted a secondary analysis to investigate how model misspecification affects the detection of causal mediation. In particular, we examined how inaccuracies in the inference model, particularly in terms of confounding variables, influence the power to detect mediation on the $T_1 - T_2$ edge. We base our simulations on $192$ Whole Blood trios that were positive for mediation under the GMAC mediation test. We consider four simulation scenarios: (i) when the inference model includes a few false positive confounders (denoted LTMS) (ii) when the inference model is missing a few true confounders with weak effects (denoted STMS) (iii) when the inference model is correct (denoted CMS) (iv) when the inference model is missing many true confounders with moderate to strong effects (SIMS)(See \autoref{txt:supptext_document} for details and methods).

Our simulations show that when the inference model contains most of the influential confounding variables, the correct inference for the $T_1 - T_2$ edge can typically be recovered (\autoref{tab:S10_Model_miss_Type_II_error_rates}). The first two simulation scenarios represent relatively minor changes in the confounding variables used in the inference model, either by omitting or including weak confounders. In these scenarios, correct inference could typically be achieved, with the exception of a small subset of trios where even the correct model failed to yield accurate results.

The final scenario (SIMS) reflects the most significant discrepancy between the true model that generates the data and model used to infer the causal model. Here, the inference model omitted an average of $28$ confounding variables (SD $\approx 7$), many of which had moderate to strong confounding effects. In this scenario, the type II error rate climbed to as high as $30\%$, underscoring the critical need to retain high-impact confounders for accurate causal inference.

\section{Discussion}
In summary, we have developed MRGNtrio to improve the power and accuracy of detecting key causal relationships in diverse genomic trios, particularly in the presence of multiple confounding variables. Our method adopts an extended interpretation of the PMR \cite{badsha2019learning} and leverages conditional and marginal hypothesis testing to identify five basic causal trio structures. By coupling MRGNtrio with a relaxed version of the adaptive confounder selection procedure described by \cite{yang2017identifying}, our method achieves efficient and precise inference of causal edges without the computational penalty typically associated with high-dimensional confounding.

Our simulations demonstrate that MRGNtrio outperforms comparable methods in detecting key causal edges ($T_1 - T_2$) and mediation models (M1) under realistic settings. A major advantage of our approach is the avoidance of the multiple testing burden; by adjusting for selected confounders within a regression framework, MRGNtrio incorporates numerous confounders without the need to infer a graph larger than three nodes. Furthermore, by integrating a permutation test for the $T_1 - T_2$ edge, our method achieves results comparable to GMAC. Our comparisons suggest that the permutation test offers a conservative measure, reducing the risk of Type I errors in the presence of non-normal data or rare genetic effects.

Our results also underscore the critical importance of accurate confounder selection in causal inference. Omitting key confounders drastically reduces the power to detect mediation, leading to Type II errors and the loss of regulatory insights. Conversely, the inclusion of false-positive confounders has a negligible impact, particularly when their effects are weak. This suggests that in practical applications, a liberal strategy favoring the inclusion of many potential confounders is preferable. Ultimately, MRGNtrio balances this requirement without excessive computational cost, making it a robust tool for genomic research.

\subsection{Algorithmic differences in confounder handling}

PC-based methods such as MRPC initialize learning by filtering edges from a fully connected graph structure. After removing edges with weak marginal relationships, the remaining edges are tested sequentially by conditioning on an increasing neighborhood of nodes \cite{badsha2019learning, spirtes2000causation}. Because the number of tests required to determine the existence of an edge is unknown a priori, MRPC utilizes sequential FDR control \cite{tian2019addis}. Consequently, the testing burden is heavily dependent on graph connectivity, signal strength, and the total number of nodes.

Critically, MRPC treats confounding variables as additional nodes within the graph. When confounders are included, the number of conditional tests increases significantly, and only edges with strong supporting evidence survive. Thus, MRPC inherently infers a larger network than MRGNtrio for the same biological question. Additionally, the adaptive FDR can result in significance thresholds that deviate substantially from the fixed $\alpha = 0.01$ employed by MRGNtrio. As demonstrated by our simulations, this causes MRPC to act more conservatively in the presence of confounding variables.

While our simulations and real data analysis demonstrate the flexibility of MRGNtrio in robustly estimating causal relationships, the method is currently limited to trios containing a genetic variant that acts as an instrumental variable. Similar to MRPC, MRGNtrio relaxes strict instrumental variable assumptions by empirically validating the association between the variant and the phenotype(s). We are currently extending this approach to develop a general MRGN algorithm capable of accommodating larger networks and addressing other classes of confounding variables.

%\section{Discussion}
%
%In summary, we have developed MRGNtrio to improve the power and accuracy of detecting key causal relationships for a genomic trio of a wide variety, in the presence of many confounding variables. Our method adopts the extended interpretation of the PMR \cite{badsha2019learning} and leverages conditional and marginal hypothesis testing to identify five basic causal structures for a trio. In addition, by coupling MRGNtrio with a more liberal version of the adaptive confounder selection procedure of \cite{yang2017identifying}, our method is capable of efficient and precise inference of causal edges in the presence of many confounders.  
%
%Our simulations demonstrate that MRGNtrio outperforms similar methods in detecting key causal edges ($T_1 - T_2$) and models such as mediation (M1) under realistic settings such as confounder selection and different types of confounding variables. Moreover, our method avoids the multiple testing burden associated with including many confounders by adjusting for selected confounders within a regression framework. As a result, MRGNtrio is capable of incorporating numerous confounders without the need to infer a graph larger than three nodes, thereby optimizing computational time and efficiency. Furthermore, by integrating the permutation test for the $T_1 - T_2$ edge, our method achieves results similar to GMAC when inferring this specific relationship. Our extensive comparisons between the t-test and permutation-based test suggest that using the permutation test provides a conservative measure, potentially reducing the risk of incorrectly inferring causal relationships in the presence of non-normal data or when dealing with rare genetic effects.
%
%Our results also underscore the critical importance of accurate confounder selection in causal inference. Omitting key confounders can drastically reduce the power to detect mediation and other causal relationships, leading to type II errors and the potential loss of crucial regulatory insights. In contrast, including false positive confounders tends to have a minor impact, especially when they have weak effects. This suggests that in practical applications, a conservative approach favoring the inclusion of more confounders may be preferable to achieve accurate causal inference. Ultimately, MRGNtrio's balanced approach provides an effective solution for incorporating confounders without excessive computational burden, while its focus on robust causal inference makes it a valuable tool in genomic research.
%
%
%\subsection{MRGNtrio and MRPC handle confounding variables differently}
%
%PC-based methods like MRPC initialize the graph learning procedure by filtering edges from the fully connected graph structure. After filtering edges with weak marginal relationships, the remaining edges are tested sequentially by conditioning on an increasing neighborhood of nodes \cite{badsha2019learning, spirtes2000causation}. For an edge to be kept in the graph, it must remain positive across a large number of conditional tests. Since the number of tests needed to determine the existance of each edge is not known beforehand, MRPC utilizes sequential FDR control to maintain the FDR at a desired threshold \cite{tian2019addis}. Therefore, the number of tests performed to determine each edge in the graph is dependent on the connectivity of the graph, the strength of the dependencies, and the total number of nodes. 
%
%MRPC treats confounding variables as additional nodes in the graph of each trio and the number of conditional tests needed to test remaining edges post-filtering is typically larger. As a result, when confounding variables are included in the network, only edges with strong supporting evidence remain in the final inferred graph. Thus when confounders are utilized in the inference for a trio, MRPC is always inferring a larger network than MRGNtrio. Also, the adaptive FDR can lead to a different cutoff values that vary widely from the cutoff of $\alpha = 0.01$ used by MRGNtrio. One consequence is MRPC tends to act conservatively in the presence of confounding variables demonstrated by our simulation results.
%
%While our simulations and real data analysis demonstrate the power of our MRGNtrio method and its flexibility to include confounding variables for robust estimation of causal relationships, it is limited to trios with a genetic variant which can act as an instrumental variable. Similar to MRPC, MRGNtrio slightly relaxes the typical assumptions on instrumental variables by checking whether the association between this variable and one or both phenotypes is significant according to the data.  
%On the other hand, MRGNtrio still accounts only for confounders.  We are extending this approach to develop a general MRGN algorithm to accommodate larger networks with more genetic variants and phenotype nodes, as well as the other two types of confounding variables.  

\end{doublespace}

\section*{Data access}
Data from The Genotype-Tissue Expression (GTEx) Project version 8 were obtained from dbGaP (\url{http://www.ncbi.nlm.nih.gov/gap}) through 
accession number phs000424.v8.pht002741.v8.p2.  Open-source R packages used in this manuscript include \url{https://github.com/Jarred6068/MRGN}, MRGNtrio (version 1.3.7.3);  \url{https://github.com/audreyfulab/MRGN}, MRPC (version 3.1.0); \url{https://cran.r-project.org/web/packages/MRPC/index.html}, and GMAC (version 3.1; 
\url{https://cran.r-project.org/web/packages/GMAC/index.html}.  Detailed results generated by our analyses are included in the supplementary 
material.

\section*{Competing interest statement}
The authors declare no competing interests.


\section*{Acknowledgments}
We acknowledge the extensive computing support provided by the Research Computing and Data Services from the Institute for Interdisciplinary Data Sciences at the University of Idaho.

This research is supported by the National Institutes of Health (NIGMS P20GM104420).

The Genotype-Tissue Expression (GTEx) Project was supported by the Common Fund of the Office of the Director of the National Institutes of Health. Additional funds were provided by the NCI, NHGRI, NHLBI, NIDA, NIMH, and NINDS. Donors were enrolled at Biospecimen Source Sites funded by NCI/Leidos Biomedical Research, Inc.\ subcontracts to the National Disease Research Interchange (10XS170), Roswell Park Cancer Institute (10XS171), and Science Care, Inc. (X10S172). The Laboratory, Data Analysis, and Coordinating Center (LDACC) was funded through a contract (HHSN268201000029C) to the The Broad Institute, Inc. Biorepository operations were funded through a Leidos Biomedical Research, Inc.\ subcontract to Van Andel Research Institute (10ST1035). Additional data repository and project management were provided by Leidos Biomedical Research, Inc.(HHSN261200800001E). The Brain Bank was supported supplements to University of Miami grant DA006227. Statistical Methods development grants were made to the University of Geneva (MH090941 \& MH101814), the University of Chicago (MH090951,MH090937, MH101825, \& MH101820), the University of North Carolina - Chapel Hill (MH090936), North Carolina State University (MH101819), Harvard University (MH090948), Stanford University (MH101782), Washington University (MH101810), and to the University of Pennsylvania (MH101822). The datasets used for the analyses described in this manuscript were obtained from dbGaP at \url{http://www.ncbi.nlm.nih.gov/gap} through dbGaP accession number phs000424.v8.pht002741.v8.p2.

\section*{Author Contributions}
J.K and A.Q.F designed the study, J.K developed the tools and analyzed the data, and J.K and A.Q.F wrote the manuscript.


\bibliography{./references.bib} 
\bibliographystyle{unsrt}

\newpage
\section*{Figures}




\begin{figure}[htbp]
	\centering
	\includegraphics[width = \textwidth]{C:/Users/Bruin/OneDrive/Documents/GitHub/MRGN_extra/Manuscript_2025/figures/5 causal models2.pdf}
	\caption[Causal models for a genetic trio]{\textbf{Causal models for a genetic trio.}
	(A) A causal diagram of a genomic trio including a genetic variant ($V$), two molecular phenotypes ($T_1$ and $T_2$), and a set of confounding variables ($\mathbf{U}$). Regression coefficients from \autoref{eq:reg1} and \autoref{eq:reg2} correspond to the directed edges in the graph. (B) The five possible causal models for a trio under the principle of Mendelian randomization (adapted from \cite{badsha2019learning}). All other configurations are categorized as ``other."}
	\label{fig:trio}
\end{figure}





\begin{figure}[htbp]
	\centering
	\includegraphics[width=\textwidth]{"C:/Users/Bruin/OneDrive/Documents/GitHub/MRGN_extra/Manuscript_2025/supplementary_figures/SF4_Confounder_Selection_Performance.pdf"}
	\caption{\textbf{Precision–recall curves comparing confounder selection performance of the CSFDR and CSnoFDR methods.} Recall is defined as the proportion of true confounders correctly identified and selected for the inference model while precision is the proportion of true confounders selected out of all confounders selected. }
	\label{fig:SF4_Confounder_Selection_Performance}
\end{figure}


\begin{figure}[H]
	\centering
	\includegraphics[width = \textwidth]{C:/Users/Bruin/OneDrive/Documents/GitHub/MRGN_extra/Manuscript_2025/supplementary_figures/SF7_Confounder_Selection_Impact_On_Inference.pdf}
	\caption{\textbf{Relationship between confounder selection and causal inference performance measured as precision and recall of edge prediction.} (A) Edge-based performance of MRGNTrio on simulated trios as a function of the number of true confounders falsely omitted from the analysis. (B) Edge-based performance of MRGNTrio on simulated trios as a function of the number of falsely included confounders.}
	\label{fig:SF7_Confounder_Selection_Impact_On_Inference}
\end{figure}



\begin{figure}[H]
	\centering
	\includegraphics[width = \textwidth]{C:/Users/Bruin/OneDrive/Documents/GitHub/MRGN_extra/Manuscript_2025/figures/MF3_GTEx.model.and.t1t2.edge.bargraphs.pdf}
	\caption[Performance of MRGNtrio on GTEx human whole blood trios]{
		\textbf{Comparative analysis of MRGNtrio, MRPC \cite{Badsha2019} and GMAC \cite{yang2017identifying} on GTEx human whole blood trios.} 
		We applied MRGNtrio (standard), MRGNtrio with permutation testing (All Perm), GMAC, and MRPC (using ADDIS FDR control) to genomic trios. 
		(A) Count of inferred causal models (M1--M4 and Other) comparing MRGNtrio and MRPC. 
		(B) Comparison of specific mediation pathways (via cis-gene vs. trans-gene) identified by MRGNtrio and GMAC. 
		(C) Total counts of trios where the $T_1 - T_2$ causal edge was inferred as Present or Absent across all four methods[cite: 33, 34, 36]. 
		(D) Histogram displaying the frequency distribution of the number of Principal Components (PCs) selected as confounders by the CSnoFDR method.
	}
	\label{fig:main_GTEx}
\end{figure}







\newpage 
\section*{Tables}



\begin{table}[H]
	\centering
	\caption{\textbf{Conditional and marginal dependence tests used by MRGNTrio.} This table summarizes the tests applied to determine which model is supported by the data. The first three models each have two possible configurations. Coefficients $\beta_{11}$ and $\beta_{12}$ come from the regression in \autoref{eq:reg1}, and $\beta_{21}$ and $\beta_{22}$ from the regression in \autoref{eq:reg2}; these coefficients are used to assess conditional independence. Marginal independence is tested using the correlations $r_{V, T_1}$ and $r_{V, T_2}$.}
	\renewcommand{\arraystretch}{0.9}
	%\renewcommand{\arraystretch}{0.9} % Shrink space between rows
	\begin{tabular}{c|p{2.8cm}|cccc|cc} 
		\toprule[0.04cm]
		\multicolumn{2}{c|}{} & \multicolumn{4}{c}{Conditional Tests} & \multicolumn{2}{|c}{Marginal Tests}  \\ 
		\midrule
		
		\multicolumn{2}{c|}{Dependence Relation} 
		& $T_1\indep V$ 
		& \small $T_1 \indep T_2$ 
		& \small $T_2\indep V$
		& \small $T_2\indep T_1$ 
		& \small $V\indep T_1$ 
		& \small $V \indep T_2$    \\
		
		\multicolumn{2}{c|}{Conditioned on} 
		& \small $T_2, {\bf U}$
		& \small $V, {\bf U}$
		& \small $T_1, {\bf U}$
		& \small $V, {\bf U}$
		& \small -- 
		& \small --    \\
		
		\multicolumn{2}{c|}{ } & \multicolumn{4}{c|}{ } & \multicolumn{2}{c}{ }  \\
		
		Model & Configuration & $\beta_{11}$ & $\beta_{12}$ & $\beta_{21}$ & $\beta_{22}$ & $r_{V, T_1}$ &  $r_{V, T_2}$ \\ 
		\midrule[0.04cm]
		
		\multirow{4}{*}{M$_0$} & \small $V \rightarrow T_1$; $T_2$ is singleton & $\neq 0$ & $= 0$ & $=0$ & $=0$ & -- &  --  \\ 
		
		& \small $V \rightarrow T_2$; $T_1$ is singleton &  $= 0$ & $= 0$ & $\neq 0$ & $= 0$ & -- & -- \\ 
		\midrule
		
		\multirow{2}{*}{M$_1$} \small & $V \rightarrow T_1 \rightarrow T_2$ & $\neq 0$ & $\neq 0$ & $= 0$ & $\neq 0$ & -- & -- \\  
		
		& \small $V \rightarrow T_2 \rightarrow T_1$ & $ = 0$ & $\neq0$ & $\neq0$ & $\neq 0$ & -- & -- \\ 
		\midrule
		
		\multirow{2}{*}{M$_2$} & $V\rightarrow T_1 \leftarrow T_2$ & $\neq0$ & $\neq 0$ & $\neq0$ & $\neq0$ & $\neq0$ & $=0$ \\ 
		
		& \small $V\rightarrow T_2 \leftarrow T_1$ & $\neq0$ & $\neq 0$ & $\neq0$ & $\neq0$ & $=0$ & $\neq 0$ \\  
		\midrule
		\multirow{1}{*}{M$_3$} & \small $T_1 \leftarrow V \rightarrow T_2$ & $\neq 0$ & $= 0$ & $\neq0$ & $=0$ & -- & -- \\  
		\midrule
		\multirow{2}{*}{M$_4$} & \small $T_1 \leftarrow V \rightarrow T_2$; $T_1 \longleftrightarrow T_2$ & $\neq0$ & $\neq 0$ & $\neq0$ & $\neq0$ & $\neq 0$ & $\neq 0$ \\  
		\bottomrule[0.04cm]
	\end{tabular}
\end{table}





\begin{table}[H]
	\caption{\textbf{Recall and precision for MRGNtrio and MRPC-ADDIS on simulated trios with confounding variables.} The highest precision and recall for each model or edge is highlighted in each row. The first column indicates the performance of MRGNtrio when the true set of confounders is used in inference for comparison. We also consider two selection methods for inferring the confounders for each trio: (i) FDR correction with the qvalue method to control the FDR at $5\%$ (CSFDR) and (ii) no FDR correction and significance taken at the $\alpha < 0.01$ level (CSnoFDR). Additional columns represent performance for MRPC-ADDIS using the CSDFR method for selecting confounders. Rows 1-5 indicate class-based performances for each of the 5 causal models, row 6 gives performance for predicting the $T_1-T_2$ edge ($\ast$ indicates similar values were observed for GMAC), and the last row gives the performance across all edges in the network using the weighting scheme described in the Section ``Performance metrics".}
	\renewcommand{\arraystretch}{0.9} % Shrink space between rows
	\begin{tabular}{l|ll|ll|ll|ll}
		\toprule
		 & \multicolumn{2}{p{2.8cm}|}{\centering \bf MRGNtrio + GT} &  \multicolumn{2}{p{2.8cm}|}{\centering \bf MRGNtrio + CSFDR} & 
		\multicolumn{2}{p{2.8cm}|}{\centering \bf MRGNtrio + CSnoFDR} & 
		\multicolumn{2}{p{2.8cm}}{\centering \bf MRPC-ADDIS + CSFDR} \\

		Metric & Recall & Prec. & Recall & Prec. & Recall & Prec. & Recall & Prec. \\ 
		\midrule[0.04cm]
		Null Model ({\bf M0}) & 0.562  & 0.995  & 0.452  & \bf 1.000  & 0.477  & 0.951  & \bf 0.836  & 0.532  \\
		Mediation ({\bf M1}) & 0.982  & 0.516  & 0.877 & \bf 0.438  & \bf 0.883  & 0.429  & 0.097  & 0.147  \\
		V-structure ({\bf M2}) & 0.919  & 0.985  & \bf 0.892  & 0.925  & 0.842  & \bf 0.968  & 0.377  & 0.443  \\ 
		Cond. Indep. ({\bf M3}) & 0.515  & 0.989  & 0.283  & \bf 0.910  & 0.305  & 0.872  & \bf 0.609  & 0.828  \\ 
		Fully Connect. ({\bf M4}) & 0.600  & 0.566  & \bf 0.501  & 0.425  & 0.464  & 0.419  & 0.217  & \bf 0.592  \\ 
		All edges  & 0.945  & 0.893  & \bf 0.931  & 0.864  & 0.911  & 0.857  & 0.770  & \bf 0.892  \\
		$T_1-T_2$ edge & 1.000  & 0.767  & \bf 0.994$^*$  & 0.703$^*$  & 0.9787 & 0.7071 & 0.7884 & \bf 0.916 \\
		\bottomrule
		%\multicolumn{9}{c}{\footnotesize $\ast$ same values observed for GMAC}
	\end{tabular}
	\label{tab:results}
\end{table}

\begin{table}[H]
	\caption{\textbf{Inference times for MRGNtrio, MRPC-ADDIS, and GMAC} on simulated trios. MRGNtrio shows a substantially faster median inference time per trio compared to both MRPC-ADDIS and GMAC, demonstrating its computational efficiency in the presence of confounding.}
	\label{tab:time_comp}
	\centering
	\begin{tabular}{lcc}
		\toprule
		\textbf{Method} & \textbf{Median time (s)} & \textbf{IQR (s)} \\
		\midrule
		MRGNTrio & 0.0105 & 0.0052 \\
		MRPC-ADDIS & 0.4601 & 9.0117 \\
		GMAC & 1.5373 & 1.4831 \\
		\bottomrule
	\end{tabular}
\end{table}

\begin{table}[H]
	\caption{\textbf{A summary of type II error rates for MRGNtrio under model misspecification}. This table reports the type I and type II error rates under four simulation scenarios investigating model misspecification. Small True Model Simulation (STMS): The trans gene is simulated using only the most significant confounders, while the inference model includes additional, less significant confounders, representing model overspecification. Large True Model Simulation (LTMS): The trans gene is simulated with a larger set of confounders, but the inference model omits weaker confounders, representing model underspecification. Correct Model Simulation (CMS) - The same model is used for both simulation and inference, assessing the stability of the permutation test when the model is correctly specified. Small Inference Model Simulation (SIMS) - The trans gene is generated using a full model, but inference is performed with only a small subset of top confounders.}
	\label{tab:S10_Model_miss_Type_II_error_rates}
	\begin{tabular}{lcccc}
		\toprule
		\textbf{Test Type} & \textbf{STMS} & \textbf{LTMS} & \textbf{CMS} & \textbf{SIMS} \\
		\midrule
		Permutation test (Type II error) & 0.1302 & 0.1302 & 0.1250 & 0.3125 \\
		Parametric test (Type II error) & 0.0729 & 0.0729 & 0.0729 & 0.2969 \\
		\bottomrule
	\end{tabular}
\end{table}

\section*{Supplementary Material}

\newpage
\subsection*{Supplementary figures}
\begin{suppfigure}[H]
	\centering
	\includegraphics[width = \textwidth]{C:/Users/Bruin/OneDrive/Documents/GitHub/MRGN_extra/Manuscript_2025/supplementary_figures/Confounding Vars2.pdf}
	\caption{\textbf{The three types of confounding variables to consider in causal inference of genetic trios}. A confounder (left) is a gene which regulates both molecular phenotypes, an intermediate variable (center) is a gene which mediates part of the effect of the genetic variant and its target, and a common child (right) is a shared target of both molecular phenotypes.}
	\label{fig:confounding}
\end{suppfigure}


\begin{suppfigure}[H]
	\includegraphics[width = \textwidth]{C:/Users/Bruin/OneDrive/Documents/GitHub/MRGN_extra/Manuscript_2025/figures/MF1_MRGN.GMAC.MRPC.all.params.t1t2.pdf}
	\caption{\textbf{Precision-recall curves for predicting the $T_1 - T_2$ edge across multiple simulation parameters}. (A) the number of simulated hidden confounders (B) Minor allele frequency of the genetic variant, (C) The signal strength of the $T_1 - T_2$ edge, and (D) The standard deviation of the regression residuals.}
	\label{fig:main_t1_t2_inf}
\end{suppfigure}

\begin{suppfigure}[H]
	\includegraphics[width = \textwidth]{C:/Users/Bruin/OneDrive/Documents/GitHub/MRGN_extra/Manuscript_2025/supplementary_figures/SF_MRGN.MRPC.all.params.edge.based.pdf}
	\caption{\textbf{Precision and recall curves of edge-based prediction performance for MRGNtrio, and MRPC with confounder selection}. (A) the number of simulated hidden confounders (B) Minor allele frequency of the genetic variant, (C) The signal strength of the $T_1 - T_2$ edge, and (D) The standard deviation of the regression residuals.}
	\label{fig:main_all_edge}
\end{suppfigure}

\begin{suppfigure}[H]
	\centering
	\includegraphics[width=4.4in]{C:/Users/Bruin/OneDrive/Documents/GitHub/MRGN_extra/Manuscript_2025/supplementary_figures/SF1_comp_times_for_conf_selection_updated.pdf}
	\caption{\textbf{Comparison of computation times between our implementation of the CSFDR method for confounder selection and that of the GMAC method \cite{yang2017identifying}}. Computation times are shown across different numbers of input trios and sizes of the candidate confounder pool.}
	\label{fig:SF1_comp_times_for_conf_selection_updated}
\end{suppfigure}


\begin{suppfigure}[H]
	\centering
	\includegraphics[width=\textwidth]{C:/Users/Bruin/OneDrive/Documents/GitHub/MRGN_extra/Manuscript_2025/supplementary_figures/SF3_GTEx_Permutation_and_Model_Changes.pdf}
	\caption{\textbf{Impact of the causal permutation test on inference and categorization of model changes for GTEx whole blood trios.} 
		(A) A $p$-value plot ($\log_{10}$ scale) comparing the observed $p$-value from the Wald test for $\beta_{12}: T_1 \indep T_2 \mid (V,U)$ with the $p$-value from the permutation test. 
		(B) A $p$-value plot for the conditional test $\beta_{22}: T_2 \indep T_1 \mid (V,U)$ showing the permutation-based $p$-values (y-axis) versus the standard Wald test. Red dotted lines indicate the significance threshold on the log scale. 
		(C) Categorization and counts of inferred model changes resulting from using the permutation test instead of the standard Wald test. 
		(D) Minor allele frequencies of trios whose inferred models changed compared with those that did not change after applying the permutation test for $\beta_{12}$ and $\beta_{22}$.}
	
	\label{fig:SF3_GTEx_Permutation_and_Model_Changes}
\end{suppfigure}



\begin{suppfigure}[H]
	\centering
	\includegraphics[width=\textwidth]{C:/Users/Bruin/OneDrive/Documents/GitHub/MRGN_extra/Manuscript_2025/supplementary_figures/SF5_Simulation_Permutation_and_Model_Changes.pdf}
	\caption[Impact of the permutation test on causal inference for simulated trios]{Impact of the permutation test on inference and categorization of inferred model changes for simulated trios. 
		(A) $p$-value plot ($\log_{10}$ scale) comparing the observed $p$-value from the Wald test for $\beta_{12}: T_1 \indep T_2 \mid (V,U)$ with the $p$-value from the permutation test. 
		(B) Same as (A) but for the conditional test $\beta_{22}: T_2 \indep T_1 \mid (V,U)$. 
		(C) Categorization and counts of inferred model changes resulting from using the permutation test instead of the standard Wald test. 
		(D) Distribution of allele frequencies for trios whose inferred models changed using the permutation test compared with trios with no model change.}
	
	\label{fig:SF5_Simulation_Permutation_and_Model_Changes}
\end{suppfigure}


\begin{suppfigure}[H]
	\centering
	\includegraphics[width = \textwidth]{C:/Users/Bruin/OneDrive/Documents/GitHub/MRGN_extra/Manuscript_2025/supplementary_figures/SF8_Trio_Skew_Stat_Distribtuions.pdf}
	\caption{\textbf{Skewness and kurtosis for cis and trans genes by allele frequency}. Stacked histograms showing the sampling distributions of sample skewness and excess kurtosis for cis and trans genes in GTEx trios. Results are segregated by the minor allele frequency of the genetic variant, highlighting differences between common and rare alleles.}

	\label{fig:SF8_Trio_Skew_Stat_Distribtuions}
\end{suppfigure}

\begin{suppfigure}[H]
	\centering
	\includegraphics[width = \textwidth]{C:/Users/Bruin/OneDrive/Documents/GitHub/MRGN_extra/Manuscript_2025/supplementary_figures/rare_trio_distributions/trio_number_730.pdf}
	\caption{\textbf{Example of a trio with a rare variant acting as an outlier.} In this trio, a single individual homozygous for the minor allele drives non-normality in trans-gene expression relative to the population. This individual (yellow point) also shows elevated expression of the cis gene. Initially, MRGNTrio inferred model M4 for this trio, while GMAC inferred no mediation. Applying the permutation test caused MRGNTrio to recover the same inference (the null model) as GMAC.}
	\label{fig:trio_number_730}
\end{suppfigure}



\begin{suppfigure}[H]
	\centering
	\includegraphics[scale =0.5]{C:/Users/Bruin/OneDrive/Documents/GitHub/MRGN_extra/Manuscript_2025/supplementary_figures/SF9_GMAC_Valid_TypeI_error.pdf}
	\caption[Comparison of GMAC and MRGNtrio type I error rates for $T_1-T_2$ detection under confounding]{Results of simulations investigating type I error rate (false mediation) under the null model. Trios were simulated under the null model for each of seven scenarios: 2, 4, 6, 8, 10, 12 or 15 confounders, excluding common child and intermediate variables. The GMAC method was applied to trios in each scenario using CSFDR confounder selection.}
	\label{fig:SF9_GMAC_Validation_Summary}
\end{suppfigure}


\subsection*{Supplementary tables}


\begin{supptable}[H]
	\caption[Simulated trios and inference results for MRGNTrio, MRPC-ADDIS, and GMAC]{
		A complete list of simulated trios, their simulation parameters, and
		inference results under MRGNTrio, MRPC-ADDIS, and GMAC.
	}
	\label{tab:ST_all_results_simulations.csv}
\end{supptable}

\begin{supptable}[H]
	\caption[GTEx trios and inference results for MRGNTrio, MRPC-ADDIS, and GMAC]{
		All analyzed GTEx whole-blood trios, including genetic annotations
		(locus, variant ID, gene name, etc.) and inference results from
		MRGNTrio, MRPC-ADDIS, and GMAC.
	}
	\label{tab:ST_GTEx_all_trios_master.csv}
\end{supptable}



%\begin{supptable}[H]
%	\caption[MRGNtrio type I and type II error under model misspecification]{This table reports the type I and type II error rates under four simulation scenarios investigating model misspecification. Small True Model Simulation (STMS): The trans gene is simulated using only the most significant confounders, while the inference model includes additional, less significant confounders, representing model overspecification. Large True Model Simulation (LTMS): The trans gene is simulated with a larger set of confounders, but the inference model omits weaker confounders, representing model underspecification. Correct Model Simulation (CMS): The same model is used for both simulation and inference, assessing the stability of the permutation test when the model is correctly specified. Small Inference Model Simulation (SIMS): The trans gene is generated using a full model, but inference is performed with only a small subset of top confounders, allowing comparison between Kvamme et al. (2021) and the MRGNtrio method.}
%	\label{tab:S10_Model_miss_Type_II_error_rates.csv}
%\end{supptable}






\begin{supptable}[H]
	\caption[Comparison of MRGNtrio, MRPC, and GMAC on GTEx trios]{This table compares MRGNtrio, MRPC, and GMAC on results from application to GTEx whole blood tissue. Results compared between MRGNtrio with CSnoFDR confounder selection to MRPC and GMAC. }
	\label{tab:S7_GTEx_comp_MRGN_MRPC_GMAC}
\end{supptable}




\subsection*{Supplementary text}

\begin{supptext}[H]
	\caption{Detailed supplementary text to explain the four simulation scenarios presented under the section ``MRGNtrio is robust to model misspecification in causal mediation". This text also includes supplementary examples with output illustrating differences in causal inference between MRGNtrio, GMAC, and MRPC on selected GTEx and simulated trios.}
	\label{txt:supptext_document}
\end{supptext}

\end{document}



























